\subsection*{27. Statistical Significance Is Not Practical Significance}

A tiny improvement may be statistically detectable in a very large study and unimportant to the buyer. A large, important improvement may be imprecisely estimated in a small pilot. P-values do not describe the probability that the business thesis is true, and failure to cross a conventional threshold does not prove that the effect is 0.

Report \textbf{effect sizes in meaningful units}, uncertainty intervals where justified, sample sizes, design limitations, and the actual decision threshold. A 2-point score gain may matter for a narrow high-stakes competency; a 10-point gain on a poorly aligned test may not. For enterprise training, minutes saved are useful only if they translate into actual throughput, reduced errors, better service, or feasible labor reallocation.

\subsection*{28. Randomization Is a Tool, Not a Substitute for Judgment}

A randomized study can be internally persuasive and strategically irrelevant if it tests an artificial task. If the intervention improves immediate quiz scores but not the independent work the customer actually values, the founder has measured the wrong outcome with admirable rigor. Conversely, an observational case may identify a critical implementation failure even without formal causal inference.

Choose evaluation designs to match the stage and the decision. A 1st pilot may need to establish feasibility, safety, task validity, and whether the program is worth improving. A later scaled offering may need stronger comparative evidence. Do not mistake the presence of statistical machinery for a high-quality business decision.

\subsection*{29. Quasi-Experiments Can Help When Random Assignment Is Unavailable}

An organization may introduce a new training intervention in 1 team before another, use a policy threshold, or roll out a product at different times. Such settings can support quasi-experimental comparisons if assumptions are credible.

Difference-in-differences compares changes in treated and untreated groups. Its core identifying requirement is not simply that the groups look alike at baseline, but that their outcomes would have followed \textbf{parallel trends absent the intervention}, after accounting for the design's relevant complications. Interrupted time-series designs examine whether an outcome's level or trend changed after implementation; other changes at the same time remain rival explanations. Regression discontinuity can use an eligibility cutoff under suitable continuity and no-manipulation assumptions. Matching can reduce differences in observed characteristics but cannot guarantee elimination of unobserved confounding.

These methods are valuable when honestly used and dangerous when treated as decorative statistical sophistication. A small startup should not claim a causal effect merely because a spreadsheet contains a control group column.

\subsection*{30. Attrition Can Reverse the Story}

Suppose 20 learners start a course. 16 finish, and 14 of those 16 pass the final task. A founder reports \textbf{87.5\% mastery among completers}. That number is arithmetically correct: 14/16. But \textbf{70\% of all starters} demonstrated the standard: 14/20, assuming none of the 4 noncompleters have verified mastery. The difference matters.

If the 4 missing students were struggling, the completer statistic may substantially overstate the outcome experienced by a typical enrollee. If they were already competent and left for unrelated reasons, the bias may have another direction. Report \emph{both} denominators, why people were missing where known, and any conservative sensitivity bounds. Never quietly remove dissatisfied or unreachable customers from the metric that is used to sell the next cohort.

\subsection*{31. A Founder Can Use Bayesian Reasoning Without Inventing Certainty}

Small early experiments often produce sparse data. Bayesian reasoning can help make assumptions explicit: prior beliefs, likelihood of observations under alternative hypotheses, and the resulting uncertainty. It may be useful to model different possible completion, mastery, or payment rates and update after each cohort.

But prior selection matters, and a posterior interval is only as meaningful as the model and data. A precise probability drawn from 10 biased participants can be less useful than an honest qualitative diagnosis. For early-stage decisions, simple counts, transparent intervals, scenario ranges, and explicit stop conditions may be better than a complicated hierarchical model. Mathematical sophistication is valuable when it improves decisions, not when it gives weak evidence an aura of objectivity.
